\subsection{回転後の重なり}

ここまでは,縦の線分が回転軸の片側にある場合を考えた.では,線分が軸をまたぐとどうなるだろうか.

縦の線分 $g(x)\le y\le f(x)$ が回転軸 $y=c$ をまたぐと,回転後の断面には中心まで点が入る.したがって内半径は $0$ であり,外半径は両端のうち軸から遠いほうの距離である.まとめると,
\[
R(x)=\max\{|f(x)-c|,|g(x)-c|\},
\]
\[
r(x)=
\begin{cases}
0,&g(x)\le c\le f(x),\\
\min\{|f(x)-c|,|g(x)-c|\},&\text{それ以外}
\end{cases}
\]
として,断面積は $\pi\{R(x)^2-r(x)^2\}$ になる.

軸より上の部分と下の部分を別々に回すと,同じ円板の部分を2度作ることがある.2つの体積を加えると,その部分を重複して数えてしまう.先に回転後の断面を考え,できた円板や円環を合わせた図形の面積を求める必要がある.

\begin{example}{軸をまたぐ領域を回す}
$0\le x\le2$,\,$x-1\le y\le1$ で表される領域を $x$ 軸のまわりに1回転してできる立体の体積を求めよ.
\end{example}

\paragraphbf{方針}
線分が軸をまたぐかどうかが変わる $x=1$ で区間を分け,円板と円環の断面積をそれぞれ積分する.

\paragraphbf{解答}
$\frac{5\pi}{3}$

\paragraphbf{解法}
$0\le x\le1$ では $x-1\le0\le1$ なので,線分は軸をまたぐ.上端の距離は $1$,下端の距離は $1-x\le1$ だから,断面は半径 $1$ の円板になる.\,$1\le x\le2$ では線分全体が軸の上にあり,断面は外半径 $1$,内半径 $x-1$ の円環になる.図\ref{図:回転断面の区別}は,この違いを断面で示している.よって
\begin{align*}
V&=\pi\int_0^1 1\dd x+\pi\int_1^2\{1-(x-1)^2\}\dd x\\
 &=\pi+\pi\left[x-\frac{(x-1)^3}{3}\right]_1^2\\
 &=\pi+\pi\left(2-\frac13-1\right)\\
 &=\frac{5\pi}{3}.
\end{align*}

\begin{figure}[htbp]
\centering
\begin{tikzpicture}[scale=1.0,>=Stealth]
\begin{scope}[xshift=-3.6cm]
\fill[blue!12](0,-1)--(0,1)--(2,1)--cycle;
\draw[->](-0.3,0)--(2.5,0)node[right]{$x$};
\draw[->](0,-1.3)--(0,1.5)node[above]{$y$};
\draw[thick](0,1)--(2,1)--(0,-1);
\draw[blue!70!black,very thick](0.5,-0.5)--(0.5,1);
\draw[orange!85!black,very thick](1.5,0.5)--(1.5,1);
\node[below]at(1,0){$1$};\node[below]at(2,0){$2$};
\node[below]at(1,-1.7){回す前の領域};
\end{scope}
\begin{scope}[xshift=0.1cm]
\fill[blue!12](0,0)circle(1);
\draw[blue!70!black,thick](0,0)circle(1);
\draw[->](0,0)--(0.707,0.707)node[midway,left]{$1$};
\fill(0,0)circle(0.025);
\node[below]at(0,-1.45){$x=\frac12$：円板};
\end{scope}
\begin{scope}[xshift=3.2cm]
\fill[orange!17,even odd rule](0,0)circle(1)(0,0)circle(0.5);
\draw[orange!85!black,thick](0,0)circle(1);
\draw[orange!85!black,thick](0,0)circle(0.5);
\draw[->](0,0)--(0.707,0.707)node[midway,right]{$1$};
\draw[->](0,0)--(-0.5,0)node[midway,above]{$\frac12$};
\fill(0,0)circle(0.025);
\node[below]at(0,-1.45){$x=\frac32$：円環};
\end{scope}
\end{tikzpicture}
\caption{軸を含む線分は円板を作り,軸から離れた線分は円環を作る.}\label{図:回転断面の区別}
\end{figure}

最初から $\pi\int_0^2\{1-(x-1)^2\}\dd x$ とすると,\,$x<1$ でも中心の半径 $1-x$ の円板をくり抜いてしまう.しかし,元の線分は軸を含んでいるので,回転後にその穴はない.

1つの $x$ に対して元の領域が複数の線分に分かれる場合には,各線分が回転して覆う半径の範囲を調べる.範囲が重なるなら,その部分を1度だけ数える.元の平面図形では重なっていない2つの部分を回す場合にも,この確認が必要になる.
